\documentclass[11pt,a4paper]{article}

% !TEX program = tectonic
% English translation of the LaTeX transcription of the three supplied handwritten pages.
% Editorial clarifications: n is the general exponent, n=2m in the even case;
% the multiplication rules for i,j and the domain of the polar form are explicit.
% The conjectures in the source are retained as conjectures.

\usepackage[T1]{fontenc}
\usepackage[utf8]{inputenc}
\usepackage[english]{babel}
\usepackage{lmodern}
\usepackage{microtype}
\usepackage{amsmath,amssymb,amsthm,mathtools}
\usepackage[
  a4paper,left=21mm,right=21mm,top=22mm,bottom=20mm,
  headheight=14pt,headsep=7mm,footskip=10mm
]{geometry}
\usepackage{fancyhdr}
\usepackage{titlesec}
\usepackage[hidelinks]{hyperref}

\hypersetup{
  pdftitle={Real Parts of Powers of z = x + iy + jt},
  pdfsubject={English LaTeX version of three handwritten mathematical notes},
  pdfauthor={},
  pdfkeywords={Powers, Real part, Polar form, Laplace-Beltrami, Conjecture}
}

\newcommand{\ii}{\mathrm{i}}
\newcommand{\jj}{\mathrm{j}}
\newcommand{\I}{\mathrm{I}}
\DeclareMathOperator{\Real}{Re}
\newcommand{\R}{\mathbb{R}}
\newcommand{\N}{\mathbb{N}}
\DeclarePairedDelimiter{\abs}{\lvert}{\rvert}

\pagestyle{fancy}
\fancyhf{}
\fancyhead[L]{\small\textsc{Mathematical Notes}}
\fancyhead[R]{\small $z=x+\ii y+\jj t$}
\fancyfoot[L]{\small Real Parts of Powers}
\fancyfoot[R]{\small\thepage}
\renewcommand{\headrulewidth}{0.3pt}
\renewcommand{\footrulewidth}{0pt}

\setlength{\parindent}{0pt}
\setlength{\parskip}{5pt}
\setlength{\jot}{4pt}
\setlength{\abovedisplayskip}{8pt plus 2pt minus 2pt}
\setlength{\belowdisplayskip}{8pt plus 2pt minus 2pt}
\setlength{\abovedisplayshortskip}{4pt plus 2pt}
\setlength{\belowdisplayshortskip}{6pt plus 2pt minus 2pt}
\titleformat{\section}{\Large\bfseries}{}{0pt}{}
\titleformat{\subsection}{\normalsize\bfseries}{}{0pt}{}
\titlespacing*{\section}{0pt}{0pt}{11pt}
\titlespacing*{\subsection}{0pt}{12pt}{5pt}
\raggedbottom

\begin{document}

% -----------------------------------------------------------------------------
% Page 1: even powers and the conjectures in the source.
% -----------------------------------------------------------------------------
\section*{Real Parts of the Powers \texorpdfstring{$z^n$}{z to the power n}}

Let
\[
  z=x+\ii y+\jj t,\qquad (x,y,t)\in\R^3,
\]
where the quaternionic multiplication rules
$\ii^2=\jj^2=-1$ and $\ii\jj=-\jj\ii$ are assumed.
For even exponents, write $n=2m$, $m\in\N$, $m\geq 1$.

\subsection*{Examples of even powers}
\begin{align*}
  \Real(z^2)
    &=x^2-(y^2+t^2),\\
  \Real(z^4)
    &=x^4-6x^2(y^2+t^2)+(y^2+t^2)^2,\\
  \Real(z^6)
    &=x^6-15x^4(y^2+t^2)+15x^2(y^2+t^2)^2-(y^2+t^2)^3,\\
  \Real(z^8)
    &=x^8-28x^6(y^2+t^2)+70x^4(y^2+t^2)^2\\
    &\qquad-28x^2(y^2+t^2)^3+(y^2+t^2)^4,\\
  \Real(z^{10})
    &=x^{10}-45x^8(y^2+t^2)+210x^6(y^2+t^2)^2\\
    &\qquad-210x^4(y^2+t^2)^3+45x^2(y^2+t^2)^4-(y^2+t^2)^5.
\end{align*}

\subsection*{Conjecture concerning the Laplace--Beltrami equation}
The real parts under consideration, $u_m(x,y,t)=\Real(z^{2m})$, have the form
\[
  u_m(x,y,t)
  =\sum_{k=0}^{m}a_k\,x^{2(m-k)}(y^2+t^2)^k,
  \qquad a_0=1,
\]
with real coefficients $a_k$ that depend on $m$.
It is conjectured that these polynomials satisfy
\[
  \boxed{\;t\,\Delta u_m=\frac{\partial u_m}{\partial t}\;},
  \qquad
  \Delta=
  \frac{\partial^2}{\partial x^2}
  +\frac{\partial^2}{\partial y^2}
  +\frac{\partial^2}{\partial t^2}
\]
where $\Delta$ denotes the Laplace operator.

\subsection*{Conjectured coefficient relations}
It is further conjectured that
\[
\begin{alignedat}{3}
  a_m     &=(-1)^m,     &\qquad& m\geq 1,\\
  a_{m-1} &=(-1)^m a_1, &      & m\geq 3,\\
  a_{m-3} &=(-1)^m a_3, &      & m\geq 5,\\
  a_{m-5} &=(-1)^m a_5, &      & m\geq 7,
\end{alignedat}
\]
and, in general,
\[
  \boxed{\;
  a_{m-(2k-1)}=(-1)^m a_{2k-1},
  \qquad k\geq 1,\quad m\geq 2k+1.
  \;}
\]

\newpage

% -----------------------------------------------------------------------------
% Page 2: polar form and induction.
% -----------------------------------------------------------------------------
\section*{Computing the Power \texorpdfstring{$z^n$}{z to the power n}}

\subsection*{Polar form}
For $\rho:=\sqrt{y^2+t^2}>0$, choose $\varphi\in(0,\pi)$ such that
\[
  \abs{z}^2=x^2+y^2+t^2,
  \qquad
  \cot\varphi=\frac{x}{\sqrt{y^2+t^2}},
  \qquad
  \I=\frac{\ii y+\jj t}{\sqrt{y^2+t^2}}.
\]
The polar form of $z$ is then
\[
  z=\abs{z}\bigl(\cos\varphi+\I\sin\varphi\bigr).
\]
The multiplication rules for $\ii$ and $\jj$ imply
\[
  \I^2
  =\frac{(\ii y+\jj t)^2}{y^2+t^2}
  =\frac{-y^2+(\ii\jj+\jj\ii)yt-t^2}{y^2+t^2}
  =-1.
\]
Thus, for every integer $n\geq 1$, we have
\begin{equation}
  \boxed{\;
    z^n=\abs{z}^{n}\bigl(\cos(n\varphi)+\I\sin(n\varphi)\bigr).
  \;}
  \tag{$*$}\label{eq:power}
\end{equation}

\subsection*{Proof by mathematical induction}
\textit{Base case.}
For $n=1$, \eqref{eq:power} is precisely the polar form of $z$.

\textit{Inductive step.}
Suppose that \eqref{eq:power} holds for some $n\geq 1$. Then
\begin{align*}
  z^{n+1}
    &=z\,z^n\\
    &=z\,\abs{z}^{n}\bigl(\cos(n\varphi)+\I\sin(n\varphi)\bigr)\\
    &=\abs{z}^{n+1}
      \bigl(\cos\varphi+\I\sin\varphi\bigr)
      \bigl(\cos(n\varphi)+\I\sin(n\varphi)\bigr)\\
    &=\abs{z}^{n+1}\Bigl[
      \cos\varphi\cos(n\varphi)
      +\I\cos\varphi\sin(n\varphi)\\
    &\hspace{30mm}
      +\I\sin\varphi\cos(n\varphi)
      +\underbrace{\I^2}_{=-1}\sin\varphi\sin(n\varphi)
      \Bigr]\\
    &=\abs{z}^{n+1}\Bigl[
      \underbrace{\cos\varphi\cos(n\varphi)
      -\sin\varphi\sin(n\varphi)}_{\cos((n+1)\varphi)}\\
    &\hspace{30mm}
      +\I\underbrace{\bigl(\cos\varphi\sin(n\varphi)
      +\sin\varphi\cos(n\varphi)\bigr)}_{\sin((n+1)\varphi)}
      \Bigr]\\
    &=\abs{z}^{n+1}
      \bigl(\cos((n+1)\varphi)+\I\sin((n+1)\varphi)\bigr).
\end{align*}
This is \eqref{eq:power} with $n+1$ in place of $n$.\hfill$\square$

\subsection*{Special case on the real axis}
For $y=t=0$, the definition of $\I$ above does not apply.
In this case, we immediately have $z=x$ and $z^n=x^n$.
The polynomial formulas for the real parts therefore remain valid there as well.

\newpage

% -----------------------------------------------------------------------------
% Page 3: trigonometric identities and the examples in the source.
% -----------------------------------------------------------------------------
\section*{Trigonometric Representation of the Real Parts}

\subsection*{Remark on the polar form}
From $\cot\varphi=x/\sqrt{y^2+t^2}$ and $\varphi\in(0,\pi)$, it follows that
\begin{align*}
  \sin\varphi
    &=\frac{1}{\sqrt{1+\cot^2\varphi}}
     =\frac{1}{\sqrt{1+\dfrac{x^2}{y^2+t^2}}}\\
    &=\frac{\sqrt{y^2+t^2}}{\sqrt{x^2+y^2+t^2}}
     =\frac{\sqrt{y^2+t^2}}{\abs{z}},\\[3pt]
  \cos\varphi
    &=\frac{\cot\varphi}{\sqrt{1+\cot^2\varphi}}
     =\frac{x}{\sqrt{y^2+t^2}\,
       \sqrt{1+\dfrac{x^2}{y^2+t^2}}}\\
    &=\frac{x}{\sqrt{x^2+y^2+t^2}}
     =\frac{x}{\abs{z}}.
\end{align*}
Hence
\[
  \boxed{\;
    \sin\varphi=\frac{\sqrt{y^2+t^2}}{\abs{z}},
    \qquad
    \cos\varphi=\frac{x}{\abs{z}}.
  \;}
\]
Since $\Real(\I)=0$, \eqref{eq:power} yields
\[
  \boxed{\;\Real(z^n)=\abs{z}^{n}\cos(n\varphi).\;}
\]

\subsection*{Examples}
\begingroup
\setlength{\jot}{3pt}
\begin{align*}
  \Real(z)
    &=\abs{z}\cos\varphi
     =\abs{z}\,\frac{x}{\abs{z}}=x,\\[3pt]
  \Real(z^2)
    &=\abs{z}^2\cos(2\varphi)
     =\abs{z}^2\bigl(\cos^2\varphi-\sin^2\varphi\bigr)\\
    &=\abs{z}^2\left(\frac{x^2}{\abs{z}^2}
       -\frac{y^2+t^2}{\abs{z}^2}\right)
     =x^2-(y^2+t^2),\\[3pt]
  \Real(z^3)
    &=\abs{z}^3\cos(3\varphi)
     =\abs{z}^3\bigl(4\cos^3\varphi-3\cos\varphi\bigr)\\
    &=\abs{z}^3\left(4\frac{x^3}{\abs{z}^3}
       -3\frac{x}{\abs{z}}\right)
     =4x^3-3x\abs{z}^2\\
    &=4x^3-3x(x^2+y^2+t^2)
     =x^3-3x(y^2+t^2),\\[3pt]
  \Real(z^4)
    &=\abs{z}^4\cos(4\varphi)
     =\abs{z}^4\bigl(8\cos^4\varphi-8\cos^2\varphi+1\bigr)\\
    &=\abs{z}^4\left(8\frac{x^4}{\abs{z}^4}
       -8\frac{x^2}{\abs{z}^2}+1\right)
     =8x^4-8x^2\abs{z}^2+\abs{z}^4\\
    &=8x^4-8x^2(x^2+y^2+t^2)+(x^2+y^2+t^2)^2\\
    &=8x^4-8x^4-8x^2y^2-8x^2t^2+x^4+y^4+t^4\\
    &\qquad+2x^2y^2+2y^2t^2+2x^2t^2\\
    &=x^4-6x^2y^2-6x^2t^2+y^4+2y^2t^2+t^4\\
    &=x^4-6x^2(y^2+t^2)+(y^2+t^2)^2.
\end{align*}
\endgroup

\end{document}
